\section{Mapping the Team to the Field}
\label{sec:map}

Newell's third and fourth questions ask where a team could work and
what important issues live there.
For a working list of "where", we borrow the nine areas of the ICSE'27
research track: AI for SE; analytics; architecture and design;
dependability and security; evolution; human and social aspects;
requirements and modeling; SE for AI; and testing and analysis.

Rather than tabulate those areas separately, we score each area on the
same nine constructs as the people and add them to
Figure~\ref{fig:focus} as reference elements (the italic columns).
Each such column is an archetype: how the typical highly cited
researcher of that area would rate, judged from the track descriptions
rather than measured from a corpus.
\note{claude}{upgrade path: build each area column from its ICSE
distinguished papers of the last five years, coded like the author
papers; then people and areas share one provenance} The red tree then
answers a question a topic list cannot: who stands nearest which area.
Some pairings are strong: Menzies sits at 92\% match to analytics, Di
Nucci at 89\% to evolution, Schmid at 89\% to both requirements and
architecture, Esposito at 86\% to SE for AI.
Others are loose: Lenarduzzi's nearest areas (evolution, architecture)
reach only 72\%, and Armenti's (human and social aspects) 69\%, so
those two are less typecast than the rest of us.
Hence question three answers itself: the team's natural territories
are the areas already adjacent to one of its members.

The far columns matter as much as the near ones.
No author sits within 85\% of security (best: Esposito, 78\%) or
testing (best: 72\%); those two areas cling to each other, not to a
person.
That is the priced blind spot of this team: papers there would be
written against the grain, or with a new collaborator who brings that
column closer.
